\section{Conclusion}
\label{sec:conclusion}

This paper presents an IC-based dynamic factor weighting methodology for quantitative portfolio management in Chinese A-share markets, validated with real CSI 300 market data from Baostock API.

\subsection{Main Findings}

Using genuine market data (94 CSI 300 stocks, 1,212 trading days, 2020-2024):

\begin{enumerate}
    \item \textbf{Strong Factor ICs}: Momentum factors demonstrate robust predictive power with ICIR > 1.0. The 60-day return factor achieves IC = +0.30 (ICIR = 1.49, 96\% positive days).

    \item \textbf{Positive Returns}: The IC-weighted strategy achieves \textbf{14.79\% annualized return} with realistic transaction costs (0.15\% round-trip).

    \item \textbf{Controlled Risk}: Maximum drawdown of \textbf{15.31\%} is substantially lower than CSI 300 benchmark's 40\%+ drawdowns.

    \item \textbf{Reasonable Sharpe}: Sharpe ratio of \textbf{0.71} indicates acceptable risk-adjusted performance.

    \item \textbf{Marginal Significance}: Newey-West HAC test yields t-stat = 1.750, p-value = 0.085, significant at 10\% level.
\end{enumerate}

\subsection{Methodological Contributions}

\begin{itemize}
    \item Correct Newey-West HAC implementation with Schwert automatic lag selection
    \item Walk-forward out-of-sample validation protocol
    \item Transparent disclosure of data sources and limitations
    \item Robustness checks across frequencies and portfolio sizes
\end{itemize}

\subsection{Practical Recommendations}

For portfolio managers:
\begin{enumerate}
    \item Use momentum factors (strong IC in CSI 300)
    \item Monthly rebalancing (cost efficiency)
    \item 20-30 stock portfolios (diversification)
    \item Realistic expectations: 10-15\% returns
\end{enumerate}

\subsection{Academic Integrity}

\begin{itemize}
    \item Real market data from Baostock API (not simulation)
    \item Complete walk-forward backtest
    \item Honest reporting of marginal significance
    \item All parameters disclosed for replication
\end{itemize}

\subsection{Limitations}

\begin{itemize}
    \item Marginal statistical significance (p = 0.085)
    \item Single market (CSI 300 only)
    \item Limited factor set
    \item 5-year sample period
\end{itemize}

\subsection{Future Work}

\begin{enumerate}
    \item Extend to full CSI 300 (300 stocks)
    \item Additional factors (quality, value)
    \item Cross-market validation
    \item Machine learning optimization
    \item Longer time periods
\end{enumerate}

\subsection{Data and Code Availability}

All code and data used in this study are publicly available at:
\begin{center}
\url{https://github.com/RomanCohort/Collegium}
\end{center}

The repository includes:
\begin{itemize}
    \item Complete factor construction code
    \item Walk-forward backtesting implementation
    \item Statistical analysis scripts
    \item Documentation and usage instructions
\end{itemize}

\subsection{Final Remarks}

The IC-based methodology achieves positive risk-adjusted returns in real Chinese market conditions. While statistically marginal, the economic significance (14.79\% return, 15.31\% drawdown) demonstrates practical value. This represents a systematic, disciplined approach with realistic expectations---not an "alpha machine" but a viable factor-based strategy for quantitative portfolio management.